\section{Restricting Attacker Observations}
\label{sec:attackerobs}
To achieve RMI, we first restrict the attacker's observations using microarchitectural isolation to satisfy \sw{shm}. We describe our baseline microarchitecture, and how it enforces strong isolation between security domains with shared memory disabled.
From there, we show how we have enabled shared memory while carefully constraining what an attacker can observe to (speculative) accesses to shared memory i.e. \sw{shm}.

\subsection{Strong Microarchitectural Isolation}
\label{sec:hardwareisolation}

\begin{figure}[t]
\centering
\includegraphics[width=\columnwidth]{figures/partitioning_mem.pdf}
\caption{Resource Partitioning Between Security Domains}
\label{fig:partitioning_all}
\centering
\includegraphics[width=\columnwidth]{figures/context_switch.pdf}
\caption{Context Switching Between Security Domains}
\label{fig:context_switch}
\centering
\includegraphics[width=\columnwidth]{figures/dual_page_table.pdf}
\caption{The Dual Page Table Mechanism }
\label{fig:dual_pt}
\end{figure}

\noindent\textbf{Memory Isolation}
Physical Memory Protection (PMP) \cite{waterman2015risc} is the official RISC-V mechanism to partition physical memory and prevent unauthorized access between security domains (OS, enclaves, and SM).
However, it relies on accessing a series of registers, making it difficult to identify memory regions outside of the cores.
This makes it difficult to integrate features such as fine-grained LLC partitioning (see Section \ref{sec:dynamiccache}).
To allow for easy determination of a memory region from a physical address, 
we revive a mechanism from the
Sanctum paper \cite{costan2016sanctum}.
At its core, the mechanism divides memory into fixed-size physical memory regions (in our case 64 regions of 32MB) allocated by the SM to different security domains.
Each core uses two 64-bit registers as memory bitmaps, storing access rights for private and shared memory of the running security domain.
Citadel adopts this mechanism, enforcing checks on all memory accesses, including speculative and page-walk operations.
Figure \ref{fig:partitioning_all} illustrates this isolation, showing that only currently running security domains can access their corresponding DRAM regions.
%\jules{18}
\\

\noindent\textbf{Microarchitectural Isolation}
We base our design on MI6 \cite{bourgeat2019mi6}, which enforces strong microarchitectural isolation through partitioning, flushing, and disabling speculation in the SM.
First, all spatially shared microarchitecture is partitioned, including the LLC, the LLC pipeline, Miss Status Handling Registers (MSHR) and other resources as illustrated in Figure \ref{fig:partitioning_all}.
In particular, the LLC is statically partitioned into 64 regions to match the 64 DRAM memory regions -- we will improve this scheme in Section \ref{sec:dynamiccache}.
We pad DRAM request timing to a fixed duration that matches the slowest request. This slightly increases latency, but masks side channels from opened rows and contention in the DRAM controller.
We also pad requests coinciding with memory refresh (periodic and independent of any secrets) to a fixed duration.
We leverage interrupt logic to interpose the SM when context switching between security domains and enforce the flushing of temporally-shared microarchitecture such as all cores' pipeline state, L1 caches, TLBs, Translation Caches and MSHRs (see Figure \ref{fig:context_switch}).
The processor does not support hyper-threading and therefore is not vulnerable to related side channels such as port contention.
The processor does not include a prefetcher or SIMD units so we are not susceptible to Spectre variants exploiting them.
Finally, Citadel is not susceptible to Meltdown-style attacks, as the processor does not speculate across privilege levels (switching to a different privilege mode flushes the re-order buffer, making sure it never contains two instructions from different privilege levels).
The exhaustive list of side channels and defense mechanisms covered by MI6 can be found in \cite{bourgeat2019mi6}.

\noindent\textbf{Disabling Speculation in the SM}
The SM executes in machine mode (see Figure \ref{fig:context_switch}) which gives it access to all physical memory.
Following MI6 \cite{bourgeat2019mi6}, speculation in the SM is disabled to prevent TES attacks. 

\subsection{Enabling Shared Memory}
\label{sec:vmanddualpt}

Existing mechanisms like MMIO or non-cacheable memory~\cite{schwarz2020context, costan2016intel} are insufficient to prevent TES attacks.
They prevent some speculative memory accesses, but rely on physical addresses or page table entry bits, requiring speculative accesses to page table entries that leak address information.
Instead, we decide to adopt Sanctum's dual-page-table mechanism (Figure \ref{fig:dual_pt}) as it makes it possible to differentiate between private and shared memory based solely on virtual addresses.
Two new registers define the enclave's private virtual memory range.
When setting up the enclave, the SM ensures that physical memory mapped as private is owned by the corresponding security domain.
Memory access translation uses one of the two page tables:
\emph{Enclave page table}: For virtual addresses within the enclave memory range. The SM ensures that this table is stored in enclave private memory.
\emph{Shared page table}: For virtual addresses outside the private range (shared memory). This table is stored in shared memory.
\noindent During address translation, additional bound checks ensure that all memory addresses that are accessed belong to the same security domain (enclave or shared memory) as the translated address.
This mechanism can also be used to enable private shared memory between enclaves (not visible to the OS).
The SM simply needs to ensure that the root page table and shared memory region(s) are not accessible to the OS.
If shared memory with the OS is also required, some page tables can be located in the OS memory.

\subsection{Disentangling Private and Shared State}
\label{sec:disentangle}
We now need to carefully constrain what the attacker can observe on (speculative) accesses to shared memory.
As a design principle, read-only state used by both shared and private memory needs to be tagged to prevent aliasing while
core-private microarchitectural state entangled with adversary state through a coherence protocol need to be partitioned.
We highlight examples relevant to our design.

\noindent\textbf{Tagged ATC and TLB for Correct Address Translations} 
Both the Address Translation Cache (ATC) and Translation Lookaside Buffer (TLB) are susceptible to cross-domain aliasing.
The ATC, which caches intermediate page walker results, relies on a few upper virtual address bits, while TLB aliasing can occur with large pages if the private virtual memory range is smaller than the page granularity.
To prevent mistranslations, we tag ATC and TLB entries to indicate their association with private or shared memory.

\noindent\textbf{Speculative Transmitters in the L1-D Cache}
\label{sec:l1bypass}
While shared memory cannot be accessed speculatively, private memory can, potentially leading to speculative eviction of shared memory.
Unfortunately, the cache coherency protocol makes some private L1 microarchitectural state observable to an attacker, as writing to a memory location cached by another core takes longer than writing to an uncached location, effectively creating a side channel.
This vulnerability enables Spectre-style attacks through coherency in the L1-D cache.
In our case, partitioning the L1 would significantly impact overall system performance, including non-protected programs, due to the cache's latency sensitivity.
Instead, we opted to bypass the L1-D cache for enclave accesses to shared memory.
Because LLC is only reconfigured when no enclaves are running and the L1-D is flushed on context switch, we avoid issues related to cache coherency.
This decision is further justified by the limited temporal reuse of shared memory accesses and the observed impact on performance (see Section \ref{sec:evalsharedmemory}).
\\

\noindent\textbf{Security Analysis}
As we tagged read-only shared resources and the L1 is the only core-private state entangled with adversary state through a coherence protocol, we do not need to worry about other core-private structures.
Consequently, the attacker is only able to observe state of the shared LLC.
This effectively limits the attacker's observations to the trace of accesses, and addresses of accesses, to shared memory i.e.~\sw{shm}.